\chapter{Réalisations}
\label{ch:realisations}

Ce chapitre, cœur du dossier, présente les réalisations qui mettent en œuvre les compétences :
l'architecture logicielle, la gestion de projet, la conception (maquettes, modèle de données,
diagrammes UML), le développement (interfaces, composants métier, accès aux données), la
sécurité, puis les tests et la veille.

\section{Architecture logicielle multicouche}
\label{sec:archi}

\projet{} suit une architecture \textbf{MVC multicouche}, où chaque couche porte une
responsabilité unique \emph{et} un rôle de sécurité. La figure~\ref{fig:archi} en donne la vue
d'ensemble.

\begin{figure}[h]
\centering
% Architecture logicielle multicouche de Solage
\begin{tikzpicture}[
  node distance=4.5mm and 0mm,
  layer/.style={draw=solageBlue!55, fill=infoBg, rounded corners=2pt,
    text width=9.2cm, align=center, minimum height=0.85cm, font=\small, inner sep=3pt},
  boot/.style={layer, fill=juryBg, draw=juryGreen!65},
  infra/.style={layer, fill=gray!10, draw=solageGray!70},
  db/.style={layer, fill=solageBlue!12, draw=solageBlue!75},
  arr/.style={-{Stealth[length=2.5mm]}, semithick, solageGray}
]
\node[layer] (nav) {\textbf{Couche cliente — Navigateur}\\{\scriptsize HTML / CSS / JS vanilla · \texttt{fetch} · rendu optimiste · \texttt{escapeHtml()}}};
\node[infra, below=of nav] (edge) {\textbf{Traefik} {\scriptsize(reverse proxy, TLS)} $\;\Rightarrow\;$ \textbf{FrankenPHP / Caddy} {\scriptsize(serveur HTTP)}};
\node[boot, below=of edge] (boot) {\textbf{public/index.php} {\scriptsize— \emph{front controller}}\\{\scriptsize autoload · \texttt{session\_start} · en-têtes CSP/nosniff · cookie HttpOnly · jeton CSRF}};
\node[layer, below=of boot] (routes) {\textbf{routes/} {\scriptsize— URL $\to$ \texttt{Controller\#method} + middleware}};
\node[layer, below=of routes] (ctrl) {\textbf{modules/controllers/} {\scriptsize— orchestration, autorisation, validation}};
\node[layer, below=of ctrl] (valid) {\textbf{modules/validators/} {\scriptsize— validation métier pure}};
\node[layer, below=of valid] (model) {\textbf{modules/models/} {\scriptsize— \textbf{tout le SQL} : PDO, requêtes préparées}};
\node[db, below=of model] (db) {\textbf{PostgreSQL} {\scriptsize— SGBD relationnel : contraintes, FK, index}};

% Vues (présentation) sur le côté
\node[layer, right=7mm of ctrl, text width=2.7cm, fill=solageBlue!8, font=\scriptsize]
  (views) {\textbf{modules/views/}\\ sortie HTML/JSON\\ \texttt{Utils::e()}};

\foreach \a/\b in {nav/edge, edge/boot, boot/routes, routes/ctrl, ctrl/valid, valid/model, model/db}
  \draw[arr] (\a) -- (\b);
\draw[arr, dashed] (ctrl.east) -- (views.west);

% étiquette src/ (framework) à gauche
\node[rotate=90, anchor=south, font=\scriptsize\itshape, text=solageGray]
  at ($(routes.west)!0.5!(model.west) + (-5mm,0)$)
  {src/ : Router · Middlewares · Logger · Utils · SessionManager · Migrations};
\end{tikzpicture}

\caption{Architecture multicouche de \projet{} : du navigateur au SGBD.}
\label{fig:archi}
\end{figure}

\subsection{Règles de séparation des couches}

Le découpage est imposé par des règles strictes que je me suis fixées, vérifiables à la lecture du code :

\begin{itemize}
  \item \textbf{aucun SQL hors de \file{modules/models/}} : un contrôleur appelle une méthode de
        modèle, il n'écrit jamais de requête ;
  \item \textbf{aucune sortie (\texttt{echo}) hors de \file{modules/views/}} et de l'entrée
        \file{public/index.php} ;
  \item \textbf{aucun accès à \texttt{\$\_POST} / \texttt{\$\_GET} / \texttt{\$\_SESSION} hors
        des contrôleurs et de \texttt{SessionManager}} ;
  \item \textbf{aucun appel direct à \texttt{Database::getConnection()} hors des modèles}, de
        \texttt{Migrations} et du \emph{bootstrap}.
\end{itemize}

Le dossier \file{src/} regroupe le code « framework » (le routeur, les middlewares, le logger,
les utilitaires, le gestionnaire de session, les migrations), distinct du code « métier »
(\file{modules/}). Cette séparation entre infrastructure réutilisable et domaine applicatif est
un marqueur d'architecture mûre.

\subsection{Rôle de sécurité de chaque couche (DICP)}

La stratégie de sécurité se lit couche par couche selon les quatre indicateurs de l'ANSSI :
\textbf{D}isponibilité, \textbf{I}ntégrité, \textbf{C}onfidentialité, \textbf{P}reuve.

\begin{table}[h]
\centering
\renewcommand{\arraystretch}{1.3}
\small
\begin{tabularx}{\textwidth}{@{}L{3cm} X X@{}}
\toprule
\textbf{Couche} & \textbf{Mécanisme de sécurité} & \textbf{Indicateur DICP} \\
\midrule
Edge (Traefik) & TLS / HTTPS, HSTS, isolation réseau & Intégrité, Confidentialité \\
Bootstrap (\file{index.php}) & En-têtes CSP / nosniff, cookie \texttt{HttpOnly} / \texttt{Secure} / \texttt{SameSite} & Intégrité, Confidentialité \\
Middlewares (\file{src/}) & Jeton CSRF, authentification, autorisation ; refus journalisé & Intégrité, Confidentialité, Preuve \\
Contrôleurs & Validation des entrées, contrôle d'accès par ressource (anti-IDOR) & Intégrité, Confidentialité \\
Modèles (PDO) & Requêtes préparées (anti-injection SQL), transactions & Intégrité \\
Vues & Échappement systématique \texttt{Utils::e()} (anti-XSS) & Intégrité \\
SGBD (PostgreSQL) & Contraintes, clés étrangères, index, comptes \& droits & Disponibilité, Intégrité \\
\bottomrule
\end{tabularx}
\caption{Rôle de sécurité de chaque couche, selon les indicateurs DICP.}
\label{tab:dicp}
\end{table}

\subsection{Patrons mis en œuvre}

L'architecture mobilise plusieurs patrons de conception : \emph{Front Controller}
(\file{public/index.php}, point d'entrée unique), \emph{MVC} (séparation modèle / vue /
contrôleur), \emph{Middleware} (chaîne de traitement transverse pour l'authentification,
l'autorisation et le CSRF) et \emph{Injection de dépendance} (le \texttt{SessionManager} reçoit
son \texttt{UserModel} par constructeur, ce qui le rend testable).

\section{Gestion de projet}
\label{sec:gestion}

Le projet a été mené en autonomie selon une démarche \textbf{itérative légère} (proche de
Kanban), structurée en phases successives : hygiène et sécurisation, conception, développement,
tests, déploiement.

\paragraph{Planification et suivi.} La feuille de route du projet (\file{ROADMAP\_DETAILLEE.md})
décrit les phases, leur état d'avancement et les priorités au regard du risque d'échec par CCP.
Git matérialise le suivi : chaque fonctionnalité correspond à une série de commits traçables, ce
qui permet de rapprocher le réalisé du planifié.

\paragraph{Comptes rendus et décisions.} Un journal de décisions techniques
(\file{Probleme-Solution.md}) consigne, pour chaque choix structurant, le contexte, le problème,
les options envisagées, la solution retenue et sa justification. Ce journal tient lieu de
comptes rendus de session et alimente la préparation de l'entretien technique.

\begin{todo}{Gestion de projet — outil collaboratif et planning visuel}
À intégrer pour compléter la compétence C4 :
\begin{itemize}
  \item une capture d'un \textbf{outil collaboratif} (GitHub Projects / Trello) montrant le
        découpage en tâches et leur avancement ;
  \item un \textbf{planning visuel} (diagramme de Gantt ou feuille de route datée) confrontant
        le \emph{planifié} au \emph{réalisé} ;
  \item trois à quatre \textbf{comptes rendus} de sessions de travail (date, objectif, réalisé,
        reste à faire).
\end{itemize}
\end{todo}

